\documentclass[aspectratio=169]{beamer}

% ===== Tema FCFM - UAdeC ====
\usetheme{FCFM}

% ===== Paquetes =====
\usepackage[utf8]{inputenc}
\usepackage{booktabs}
\usepackage[backend=biber, style=numeric, sortcites]{biblatex}
\addbibresource{referencias.bib}

% ===== Datos de Portada ====
\title[Plantilla FCFM]{Plantilla de la Facultad de Ciencias Físico-Matemáticas}
\subtitle{Prueba de texto}
\date{\today}
\author[R. Chávez B.]{Roberto Chávez Berlanga}
\institute[UAdeC]{Universidad Autónoma de Coahuila}

% ===== Inicio de Documento ====
\begin{document}
	
	\begin{frame}
		\titlepage
	\end{frame}
	
	\section{Estructura y bloques}
	
	\begin{frame}{Los bloques hacen que todo resalte}
		\begin{block}{Bloque por defecto}
			Este es un bloque por defecto.
		\end{block}
		\begin{alertblock}{Bloque de Alerta}
			Este es un bloque de alerta.
		\end{alertblock}
		\begin{exampleblock}{Bloque de Ejemplo}
			Este es un bloque de ejemplos.
		\end{exampleblock}
	\end{frame}
	
	\begin{frame}{Listas e items}
		\begin{itemize}
			\item Primer punto principal.
			\begin{itemize}
				\item Subpunto con triángulo.
				\item Otro subpunto.
			\end{itemize}
			\item Segundo punto principal.
			\item Tercer punto principal.
		\end{itemize}
	\end{frame}
	
	\section{Matemáticas}
	
	\begin{frame}{Puedes escribir matemáticas aquí}
		\begin{equation}
			a^2 = \alpha^2 + c^2
		\end{equation}
		\begin{equation}
			\int_{0}^{\infty} e^{-x^2}\, dx = \frac{\sqrt{\pi}}{2}
		\end{equation}
	\end{frame}
	
	\begin{frame}{Una tabla sencilla}
		\begin{center}
			\begin{tabular}{lcc}
				\toprule
				Método & Tiempo & Error \\
				\midrule
				Euler        & 0.12 s & 0.05 \\
				Runge--Kutta & 0.45 s & 0.001 \\
				\bottomrule
			\end{tabular}
		\end{center}
	\end{frame}
	
	\section{Código}
	
	\begin{frame}[fragile]{Ejemplo de código en Python}
		\begin{lstlisting}[language=Python]
			def factorial(n):
			# Caso base
			if n <= 1:
			return 1
			return n * factorial(n - 1)
			
			print(factorial(5))
		\end{lstlisting}
	\end{frame}
	
	\begin{frame}[fragile]{Ejemplo de código en C}
		\begin{lstlisting}[language=C]
			#include <stdio.h>
			
			int main(void) {
				// Imprime los primeros cinco enteros
				for (int i = 0; i < 5; i++) {
					printf("Valor: %d\n", i);
				}
				return 0;
			}
		\end{lstlisting}
	\end{frame}
	
	\begin{frame}[fragile]{Comentarios con acentos }
		\begin{lstlisting}[language=Python]
			# Función que calcula el área de un círculo
			def area_circulo(radio):
			# El número pi está aproximado
			pi = 3.1416
			return pi * radio ** 2  # ¿Está bien así?
		\end{lstlisting}
	\end{frame}
	
	\begin{frame}{Un tema con citas}
		Como se demostró en trabajos previos \cite{einstein1905},
		la relatividad cambió la física. Otros autores
		\cite{knuth1984, lamport1994} aportaron al cómputo.
	\end{frame}
	
	\begin{frame}[allowframebreaks]{Referencias}
		\printbibliography
	\end{frame}
	
\end{document}